

\addcontentsline{toc}{section}{Список литературы}
\begin{thebibliography}{99}

\bibitem{косарева2020}
Косарева Н.А., Чаунина Е.А., Ахмиева Г.Р., Петрова М.Ю. 
Влияние качества кормовых рационов на производство молока // 
Современное состояние, перспективы развития АПК и производства специализированных продуктов питания: 
материалы Международной научно-практической конференции, посвящённой юбилею Заслуженного работника 
высшей школы Российской Федерации, доктора технических наук, профессора Гавриловой Натальи Борисовны, 
24 апреля 2020 года. – Омск: Омский ГАУ, 2020. – С. 700–706.

\bibitem{ганущенко2019}
Ганущенко О. 
Структурность кормосмесей для коров // 
Животноводство России. – 2019. – № 12. – С. 59–61.

\bibitem{redkin2022}
Редькин С.В., Балгабаева Н.Р. 
Периоды лактации коровы с точки зрения нормальной физиологии // 
Инновационная наука. – 2022. – № 1-2. – С. 45–47. – 
ISSN 2410-6070.

\bibitem{rosstip2025}
Сено, солома, сенаж и силос: как выбрать оптимальное соотношение грубых кормов для рациона КРС в разные сезоны //
Росстип: информационный портал. – 2025. – 
URL: https://rosstip.ru/articles/seno-soloma-senazh-i-silos-kak-vybrat-optimalnoe-sootnoshenie-grubykh-kormov-dlya-ratsiona-krs-v-raznye-sezony/ (дата обращения: 17.02.2026).

\bibitem{rozhkova2021}
Рожкова-Тимина И.О.
Анализ рационов коров голштинской породы в сухостойный период в условиях Сахалина //
Животноводство и кормопроизводство. – 2021. – Т. 104, № 2. – С. 120-130. –
DOI: 10.33284/2658-3135-104-2-120.

\bibitem{baranova2023}
Баранова Н.С., Кирикова Т.Н., Давыдова А.С., Казаков Д.С.
Оценка биологической полноценности кормов по контрольным надоям молока и ее использование при организации кормления молочного скота //
Аграрная наука Евро-Северо-Востока. – 2023. – Т. 24, № 3. – С. 459-467. –
DOI: 10.30766/2072-9081.2023.24.3.459-467.

\bibitem{shishkina2025}
Шишкина Т.В., Зыкина Е.А., Семикова Н.М., Фомина В.В.
Молочная продуктивность коров в зависимости от сезона отела //
Нива Поволжья. – 2025. – № 1 (73). – 
DOI: 10.36461/NP.2025.73.1.020.

\bibitem{zolotarov2023}
Золотарев А., Родионова К., Химич М. и др.
Способы снижения влияния внешней среды в летний период на молочную продуктивность коров //
Scientific Horizons. – 2023. – Т. 26, № 4.

\bibitem{gusarov2024}
Гусаров И.В., Обряева О.Д.
Система нормированного кормления высокопродуктивных коров: монография /
И.В. Гусаров, О.Д. Обряева. – Вологда: ВолНЦ РАН, 2024. – 154 с.

\bibitem{poroda}
Министерство сельского хозяйства Российской Федерации.
Породы племенных сельскохозяйственных животных и птицы, распространенные в Российской Федерации: Каталог / 
под рук. Х.А. Амерханова, В.Ф. Федоренко. – 
М.: ФГНУ «Росинформагротех», 2006. – 60 с. – 
ISBN 5-7367-0569-9.

\bibitem{kozelsk2026}
Ветеринарная клиника в Козельске.
Основные болезни крупного рогатого скота. Инфекционные болезни крупного рогатого скота //
Сайт "Ветеринарная клиника в Козельске". – 2026. – 3 января. –
URL: https://kozelsk-vet40.ru/novosti-i-meropriyatiya/osnovnyie-bolezni-krupnogo-rogatogo-skota-infekczionnyie-bolezni-krupnogo-rogatogo-skota.html.

\bibitem{tsenovik}
Справочник "Ценовик".
Болезни КРС //
Сайт "Ценовик.ру". –
URL: https://www.tsenovik.ru/spravochnik/zootekhniya/krs/bolezni-krs/.

\bibitem{msd2021}
MSD Veterinary Manual.
Болезни коров //
Сайт "MSD Animal Health. Справочник владельца". – 2021. – 12 февраля. –
URL: https://ruminants.msd-animal-health.ru/disease/bolezni-korov/.

\bibitem{sciencemail2025}
Статья "Самые длинные имена: кто их носит, и зачем их дают" //
Наука Mail.ru. – 29.07.2025. –
URL: \url{https://science.mail.ru/articles/5830-samye-dlinnye-imena/} (дата обращения: 18.04.2026).

\bibitem{genoru2013}
Новость "Обладатель самого длинного имени в Сибири живет в Новосибирске" //
Международный институт генеалогических исследований (geno.ru). – 19.02.2013. –
URL: \url{https://geno.ru/about/news/6810/} (дата обращения: 18.04.2026).

\bibitem{minzdrav2021}
Приказ Министерства здравоохранения РФ от 20 декабря 2012 г. № 1175н 
«Об утверждении порядка назначения и выписывания лекарственных препаратов...» 
(с изменениями и дополнениями от 2 декабря 2013 г., 30 июня 2015 г., 31 октября 2017 г., 27 сентября 2021 г.) // 
Электронный фонд правовых и нормативно-технических документов. – 2021. – 
URL: \url{https://docs.cntd.ru/document/902393032} (дата обращения: 18.04.2026).

\end{thebibliography}
\begin{comment}
https://www.icar.org/guidelines/icar-central-health-key/ (словарь заболеваний) макс 119
\end{comment}

